\chapter{Introduction}
\label{chap:introduction}

% TODO: draft, expand once the rest of the thesis exists
\begin{introduction}
This chapter introduces the problem addressed by this thesis, states its objectives and contributions, and outlines the structure of the remaining chapters.
\end{introduction}

\section{Context and Motivation}
% TODO: digital twin systems for physical infrastructure are typically built as
% closed, purpose-specific applications, tightly coupled to one infrastructure
% type and one data backend. Motivate the need for a generalized, decoupled
% viewer.

\section{Thesis Objectives}
% TODO: state the objectives explicitly, e.g.:
% - design and implement an open-source digital twin viewer generalizable
%   across infrastructure types
% - decouple the viewer's rendering/interaction layer from any single data
%   backend
% - evaluate the resulting architecture through performance benchmarks and
%   user testing

\section{Contributions}
% TODO: list concrete contributions (the viewer itself as open-source
% software, the architectural pattern, the benchmark harness/results, the
% user study once completed)
\begin{itemize}
    \item % TODO
\end{itemize}

\section{Thesis Structure}
This document is organized into six chapters. Chapter~\ref{chap:introduction} introduces the problem and objectives. Chapter~\ref{chap:state-of-the-art} reviews related work on digital twin visualization and geospatial engines. Chapter~\ref{chap:architecture} presents the high-level architecture of the viewer. Chapter~\ref{chap:implementation} details the implementation decisions made to achieve generalization and decoupling. Chapter~\ref{chap:evaluation} evaluates the resulting system through performance benchmarks and user testing. Chapter~\ref{chap:conclusions} summarizes the contributions and discusses limitations and future work.
